% =====================================================================
%  2bat-ud3-5.tex  —  SESSIÓ 5: Asímptotes: la línia que no es traspassa
% =====================================================================
\horatitol{Sessió 5 — Asímptotes: la línia que no es traspassa}%
{Posar nom a la línia a la qual s'acosta una gràfica sense tocar-la mai: asímptotes horitzontals i verticals, i la seva lectura social.}

\subsection{Idea clau}

El «sostre» a què tendeixen els fenòmens té un nom: \textbf{asímptota}. Una asímptota
és la \textbf{línia recta} a la qual s'acosta la gràfica d'una funció sense arribar-hi
mai. N'hi ha de dos tipus, i cadascun té una lectura social diferent.

\begin{idea}
\textbf{Dos tipus d'asímptotes}
\begin{itemize}[leftmargin=1.4em, itemsep=3pt]
  \item \textbf{Horitzontal:} una recta horitzontal $y=L$ a la qual s'acosta la gràfica
        quan $x$ creix molt. Està lligada al \textbf{límit a l'infinit}
        ($\lim_{x\to\infty}f(x)=L$). \emph{Lectura social:} un \textbf{sostre de
        saturació} (la capacitat d'un servei, el valor on s'estabilitza una llista).
  \item \textbf{Vertical:} una recta vertical $x=a$ a prop de la qual la funció
        \textbf{es dispara} (cap a $+\infty$ o $-\infty$). Apareix en un valor que la
        funció \textbf{no pot prendre} (per exemple, un denominador que es fa zero).
        \emph{Lectura social:} una \textbf{situació impossible} o un punt crític.
\end{itemize}
\end{idea}

\subsection{Exemples}

\begin{exemple}[asímptota horitzontal: la llista que s'estabilitza]
La llista d'espera $f(x)=800-\dfrac{700}{x}$ s'acostava a $800$. La recta $y=800$ és
una \textbf{asímptota horitzontal}: la gràfica s'hi acosta sense tocar-la.

\begin{center}
\begin{tikzpicture}
\begin{axis}[udaxis, width=8.6cm, height=4.4cm,
    xlabel={\small $x$}, ylabel={\small $f(x)$},
    xmin=0, xmax=28, ymin=0, ymax=950, xtick={0,7,14,21,28}, ytick={0,400,800}]
  \addplot[udblue, domain=1:28, samples=100]{800-700/x};
  \addplot[udnavy, dashed, domain=0:28, samples=2]{800};
  \node[udnavy, font=\scriptsize, anchor=south east] at (axis cs:28,800) {asímptota $y=800$};
\end{axis}
\end{tikzpicture}

\figcap{Asímptota horitzontal $y=800$: el sostre de saturació de la llista.}
\end{center}
\end{exemple}

\begin{exemple}[asímptota vertical: la hipèrbola del fons de 1r]
El repartiment d'un fons $f(x)=\dfrac{\num{36000}}{x}$ té una \textbf{asímptota
vertical} a $x=0$: si el nombre de famílies s'acostés a zero, l'ajut per família es
dispararia cap a l'infinit. També té una \textbf{asímptota horitzontal} $y=0$ (l'ajut
tendeix a zero amb moltes famílies).

\begin{center}
\begin{tikzpicture}
\begin{axis}[udaxis, width=8.6cm, height=4.6cm,
    xlabel={\small famílies ($x$)}, ylabel={\small \euro/família},
    xmin=0, xmax=20, ymin=0, ymax=20000, xtick={0,5,10,15,20}, ytick={0,10000},
    yticklabel style={/pgf/number format/fixed}]
  \addplot[udblue, domain=2:20, samples=100]{36000/x};
  \addplot[udnavy, dashed] coordinates {(0,0) (0,20000)};
  \node[udnavy, font=\scriptsize, anchor=south west] at (axis cs:0.2,15000) {as.\ vertical $x=0$};
\end{axis}
\end{tikzpicture}

\figcap{Asímptota vertical $x=0$: amb molt poques famílies, l'ajut es dispara.}
\end{center}
\end{exemple}

\subsection{Exercicis resolts}

\begin{exresolt}[identificar les asímptotes]
Per a $f(x)=200-\dfrac{500}{x}$, digues quina és l'asímptota horitzontal i raona si en
té de vertical.
\solucio
\textbf{Horitzontal:} quan $x$ creix, $\dfrac{500}{x}$ s'acosta a $0$, així que $f(x)$
s'acosta a $200$: l'asímptota horitzontal és $y=200$. \textbf{Vertical:} la funció no
està definida a $x=0$ (no es pot dividir per zero); a prop de $x=0$ es dispara. Té
asímptota vertical $x=0$.
\resposta{Horitzontal $y=200$; vertical $x=0$.}
\end{exresolt}

\begin{exresolt}[interpretar socialment una asímptota vertical]
En el repartiment $f(x)=\dfrac{\num{36000}}{x}$, què significa, en el context real, la
asímptota vertical a $x=0$?
\solucio
Vol dir que, si el nombre de famílies fos pràcticament zero, l'ajut per família es
dispararia cap a quantitats enormes. És una situació \textbf{límit i poc realista}: el
model perd sentit quan $x$ s'acosta a $0$, perquè no té sentit repartir un fons entre
«cap» família. L'asímptota vertical marca aquí un \textbf{valor que la variable no pot
prendre} ($x=0$).
\resposta{Amb gairebé cap família l'ajut es dispararia; $x=0$ és un valor impossible
en el context.}
\end{exresolt}

\subsection{Exercicis per fer}

\begin{exercicis}
\begin{exlist}
  \item Per a cada funció, digues quina és l'asímptota horitzontal (mirant cap a quin
        valor tendeix quan $x$ creix):
        \begin{enumerate}[label=\alph*), leftmargin=1.6em, topsep=2pt]
          \item $f(x)=\dfrac{10}{x}$ \hfill \item $g(x)=150-\dfrac{300}{x}$ \hfill
          \item $h(x)=400+\dfrac{100}{x}$
        \end{enumerate}
  \item Quines d'aquestes funcions tenen asímptota vertical a $x=0$? Raona-ho.
        \begin{enumerate}[label=\alph*), leftmargin=1.6em, topsep=2pt]
          \item $f(x)=\dfrac{5}{x}$ \hfill \item $g(x)=2x+1$ \hfill
          \item $h(x)=10-\dfrac{3}{x}$
        \end{enumerate}
  \item Una funció d'ocupació té asímptota horitzontal $y=250$. Què representa aquest
        valor en el context d'un servei?
  \item Dibuixa a mà una gràfica que tingui asímptota horitzontal $y=100$ i que s'hi
        acosti per sota a mesura que $x$ creix.
  \item \textbf{(Lectura social)} Associa cada asímptota amb la seva interpretació:
        \begin{enumerate}[label=\alph*), leftmargin=1.6em, topsep=2pt]
          \item asímptota horitzontal $y=300$ \quad (\;) un valor que la variable no pot
                prendre
          \item asímptota vertical $x=0$ \quad (\;) un sostre de saturació del servei
        \end{enumerate}
  \item \textbf{(Síntesi)} Explica amb les teves paraules la diferència entre una
        asímptota horitzontal i una de vertical, i posa un exemple social de cada una.
\end{exlist}
\end{exercicis}

% ---------------------------------------------------------------------
\fullsolucions{Sessió 5}
\begin{solucions}
\begin{sollist}
  \item (a) $y=0$; \quad (b) $y=150$; \quad (c) $y=400$.
  \item (a) \textbf{sí} (denominador $x$ es fa zero a $x=0$); \quad (b) \textbf{no} (és
        una recta, definida per a tot $x$); \quad (c) \textbf{sí} (té $\frac{3}{x}$,
        que es dispara a prop de $x=0$).
  \item Representa el \textbf{sostre de saturació}: el servei s'estabilitza en $250$
        (per exemple, $250$ places ocupades), valor que no superarà.
  \item Esbós: una corba creixent que s'aplana acostant-se per sota a la recta
        horitzontal $y=100$, sense traspassar-la.
  \item (a) $\rightarrow$ un sostre de saturació del servei; \quad
        (b) $\rightarrow$ un valor que la variable no pot prendre.
  \item Resposta oberta. Idea: l'horitzontal ($y=L$) és el valor cap al qual
        s'estabilitza el fenomen quan $x$ creix (sostre); la vertical ($x=a$) és un
        valor prohibit de la variable on la funció es dispara. Exemples: horitzontal =
        capacitat d'un centre; vertical = repartir entre $0$ persones.
\end{sollist}
\end{solucions}
